% ===================== ALCANCE =====================
\section{Alcance y Limitaciones}
\label{sec:alcance}

\subsection{Alcance del Estudio}

Este estudio de validación hidráulica e ingeniería de procesos comprende los elementos definidos en la Tabla~\ref{tab:alcance_incluido}.

\begin{table}[htbp]
    \centering
    \caption{Elementos incluidos en el alcance del estudio.}
    \label{tab:alcance_incluido}
    \setcounter{itemcount}{0}
    \begin{tabularx}{\textwidth}{@{}NX@{}}
        \toprule
        \multicolumn{1}{c}{\textbf{Item}} & \textbf{Descripción} \\
        \midrule
        & Modelamiento hidráulico dinámico en estado estacionario para flujo incompresible en el tramo Km 16.5 de las líneas de salmuera y condensado. \\
        & Simulación y comparación cuantitativa de pérdidas de presión dinámicas para la tubería superficial original (\SI{450.00}{\meter}) y la proyectada por PHD (\SI{452.26}{\meter}). \\
        & Determinación de presiones hidrostáticas estáticas localizadas en el tramo horizontal de fondo (cota \SI{2542.75}{m.s.n.m.}). \\
        & Identificación cualitativa y planteamiento de mitigaciones para riesgos de sedimentación de sólidos y acumulación de bolsas de aire en el sifón invertido. \\
        \bottomrule
    \end{tabularx}
\end{table}

\subsection{Limitaciones del Estudio}

Las exclusiones técnicas y restricciones del estudio se detallan en la Tabla~\ref{tab:alcance_excluido}.

\begin{table}[htbp]
    \centering
    \caption{Elementos excluidos del alcance y limitaciones del estudio.}
    \label{tab:alcance_excluido}
    \setcounter{itemcount}{0}
    \begin{tabularx}{\textwidth}{@{}NX@{}}
        \toprule
        \multicolumn{1}{c}{\textbf{Item}} & \textbf{Descripción} \\
        \midrule
        & Modelamiento dinámico de transitorios hidráulicos de sobrepresión (golpe de ariete) para la línea de salmueroducto completa. \\
        & Ensayos de laboratorio para la caracterización fisicoquímica de muestras en sitio de salmuera o condensado; se emplearon propiedades estandarizadas del NIST. \\
        & Simulación hidráulica de la red de tuberías de transporte completa aguas arriba y aguas abajo del Km 16.5. \\
        & Diseño mecánico-estructural de la perforación horizontal dirigida o cálculo geotécnico de cargas de terreno sobre el tubo. \\
        \bottomrule
    \end{tabularx}
\end{table}
